\begin{afterword}
\summaryhead

The post argues that the usual way of picturing the Prisoner's Dilemma misleads its human readers. It sets out the game: each player prefers defecting against a cooperator, then mutual cooperation, then mutual defection, then cooperating against a defector. Defection therefore dominates, although both players prefer mutual cooperation to mutual defection.

The standard illustrations, two prisoners deciding whether to testify or two strangers playing once for dollars, ask the reader to pretend to be entirely selfish. The author holds that a human cannot do this. Evolution, through the iterated dilemma, has given us fairness and sympathy, so for a human the story's payoffs are not really those of the dilemma, and the instinct is to seek mutual cooperation.

A true dilemma needs an opponent we do not care about at all: a paperclip maximizer from another dimension, bargaining once over a substance that could cure billions of people. There, defecting against a cooperator really does seem better, yet both sides still prefer mutual cooperation to mutual defection. The post ends by asking the reader what to do.

\responsehead

The post's game theory is correct. Both matrices have the defining order for both players, the prison sentences are the standard ones, and the dominance argument holds. The paperclip case does what it was built to do. It removes every reason to care about the other player, including any future dealings, so that the numbers in the matrix are ones a human really holds.

The underlying point is standard, and the post does not claim otherwise: it aims its charge at the usual picture and accepts the analysis. Game theory already takes the entries of the matrix to include everything a player cares about. Robin Hanson said so in the comments, and the author replied that the complaint was about how the game is illustrated. Much of the rest has a history that supports the post. Kelley and Thibaut formalized the gap between material outcomes and the matrix a person acts on in 1978. A player who ranks mutual cooperation above exploiting a cooperator, and asks only whether the other will cooperate, is playing what game theorists call a stag hunt or assurance game. The evolutionary story is Trivers's theory of reciprocal altruism.

What the post adds is a vivid example and a claim about people, and the claim has thin support. That a human ``cannot pretend to be genuinely, truly selfish'' rests on an analogy with jury instructions. In the study the linked post reports, instructions did not reduce hindsight bias, but the post gives no evidence that they fail in the same way here. The descriptions of what we feel in laboratory games cite no study. Experiments find conditional cooperators and also a large minority of free riders, and the post's ``we'' fits the first group. These claims are hedged (``not entirely'', ``we may''), and Hanson's experience as a teacher agrees that students mostly hear the story the wrong way.

The post leaves its closing question open. In the comments the author wrote that it had not said the author would defect.

In short: a correct statement of the dilemma and a vivid example whose payoffs a human really holds; the point that payoffs must include everything a player cares about is standard game theory, and the claims about human instinct that motivate the example are offered without evidence.
\end{afterword}
